% ======================================================================
% main2.tex —— 《动机社会》重构版独立入口
% ----------------------------------------------------------------------
% 本文件是结构重构版（v2）的编译入口。结构与内容权威：
%   notes/00-architecture.md
% 原书入口 main.tex 与全部原章节文件保持零改动；本文件只引用
%   draft/** 下的新章节。
% 编译：xelatex main2.tex && xelatex main2.tex（两遍，出目录与交叉引用）
%
% 字体说明：macOS 下使用与原书 main.tex 完全相同的字体配置
% （Songti SC / Menlo / Times New Roman / Helvetica / STIX Two Math，
%  fontset=none）。在非 macOS 环境（Linux 沙箱 / CI）下自动回退到
% TeX Live 自带字体（Fandol / TeX Gyre / Latin Modern Math），
% 该回退分支在 macOS 上不会生效。
% ======================================================================
\documentclass[openany, a4paper, 12pt]{book}

% ==========================================
% 1. 基础宏包配置（自 main.tex 复制）
% ==========================================
\usepackage[UTF8,fontset=none]{ctex}     % 中文支持（手动指定字体，避免 TinyTeX mac fontset 缺失）
\IfFontExistsTF{Songti SC}{%
  % ---- macOS 正式字体配置（与原书 main.tex 完全一致）----
  \setCJKmainfont{Songti SC}
  \setCJKsansfont{Songti SC}
  \setCJKmonofont{Menlo}
  \setmainfont{Times New Roman}
  \setsansfont{Helvetica}
  \setmonofont{Menlo}
}{%
  % ---- 非 macOS 回退（沙箱/CI 编译验证用；macOS 上自动走上面的分支）----
  \setCJKmainfont[Extension=otf,BoldFont=FandolSong-Bold,ItalicFont=FandolKai-Regular]{FandolSong-Regular}
  \setCJKsansfont[Extension=otf,AutoFakeBold=2.2]{FandolHei-Regular}
  \setCJKmonofont[Extension=otf,AutoFakeBold=2.2]{FandolFang-Regular}
  \setmainfont{texgyretermes}[
    Extension=.otf,
    UprightFont=*-regular,
    BoldFont=*-bold,
    ItalicFont=*-italic,
    BoldItalicFont=*-bolditalic]
  \setsansfont{texgyreheros}[
    Extension=.otf,
    UprightFont=*-regular,
    BoldFont=*-bold,
    ItalicFont=*-italic,
    BoldItalicFont=*-bolditalic]
  \setmonofont{texgyrecursor}[
    Extension=.otf,
    UprightFont=*-regular,
    BoldFont=*-bold,
    ItalicFont=*-italic,
    BoldItalicFont=*-bolditalic]
}
\usepackage{unicode-math}
\IfFontExistsTF{STIX Two Math}{%
  \setmathfont{STIX Two Math}
}{%
  \setmathfont{latinmodern-math.otf}   % 非 macOS 回退
}
\usepackage{geometry}       % 页面设置
\geometry{left=1in, right=1in, top=1in, bottom=1in}
\usepackage{fancyhdr}       % 页眉页脚
\usepackage{hyperref}       % 超链接目录
\hypersetup{colorlinks=true, linkcolor=blue, filecolor=magenta, urlcolor=cyan}
\usepackage{subfiles}       % 子文件支持

% ==========================================
% 2. 绘图与框图工具 (自 main.tex 复制)
% ==========================================
\usepackage{tikz}
\usetikzlibrary{shapes.geometric, arrows.meta, positioning, shadows, calc, fit, trees}

\usepackage[most]{tcolorbox}
\tcbuselibrary{skins, breakable, theorems}

% 阅读笔记框
\newtcolorbox{readingnote}{
	colback=gray!5,
	colframe=gray!40,
	boxrule=0.4pt,
	arc=2mm,
	left=6pt,
	right=6pt,
	top=6pt,
	bottom=6pt
}

\newcommand{\xr}{\xrightarrow}

% ==========================================
% 数学附录 ↔ 正文 链接接口
% ==========================================
\newcommand{\mathlink}[2]{%
	\par\noindent
	\textcolor{blue}{\small$\rightarrow$ \hyperref[#1]{#2}}%
	\par
}

% ==========================================
% 低频高信息密度强调（写作工程接口）
% ==========================================
\newcommand{\lowfreq}[1]{\textbf{\textcolor{red!70!black}{#1}}}


% ==========================================
% 3. 自定义环境定义（自 main.tex 复制）
% ==========================================

% 定义【逻辑地图】环境
% 使用 tcolorbox 包装，带有标题和边框
\newenvironment{logicmap}{
	\begin{tcolorbox}[
		colback=blue!5!white,      % 背景色：极淡蓝
		colframe=blue!50!black,    % 边框色：深蓝
		title=\textbf{逻辑地图 (Logic Map)}, % 标题
		fonttitle=\bfseries,
		sharp corners=south,       % 底部直角
		boxrule=0.5mm,
		]
	\centering
	}{
	\end{tcolorbox}
}

% 定义【写作动机】环境
% 使用灰色背景强调
\newenvironment{motivation}{
	\vspace{0.5em}
	\begin{tcolorbox}[
		colback=gray!10,           % 背景色：淡灰
		colframe=gray!60,          % 边框色：深灰
		title=\textit{本章动机 (Motivation)},
		fonttitle=\bfseries,
		arc=0mm,                   % 直角边框
		leftrule=3mm               % 左侧加粗
		]
		\small \itshape % 字体变小并倾斜
	}{
	\end{tcolorbox}
	\vspace{1em}
}

% --- 2.4 数学推导盒子 ---
\newtcolorbox{mathbox}[1]{
	enhanced,
	colback=gray!5,
	colframe=gray!60!black,
	title=\textbf{数学建模：#1},
	fonttitle=\small\bfseries\sffamily,
	leftrule=2mm,
	breakable
}

\newtcolorbox{definitionbox}{
	colback=blue!3!white,
	colframe=blue!60!black,
	fonttitle=\bfseries,
	breakable,
	enhanced,
	boxrule=0.8pt,
	arc=2mm
}
\newtcolorbox{questionbox}{
	colback=red!3!white,
	colframe=red!60!black,
	breakable,
	enhanced,
	boxrule=0.8pt,
	arc=2mm
}

\newtcolorbox{examplebox}{
	colback=green!3!white,
	colframe=green!50!black,
	breakable,
	enhanced,
	boxrule=0.8pt,
	arc=2mm
}

\newenvironment{reader_anticipation}{
	\begin{tcolorbox}[colback=yellow!3!white,colframe=yellow!50!black,title=\textit{读者可能还会问},breakable]
	\begin{itemize}
}{
	\end{itemize}
	\end{tcolorbox}
}

\usepackage{booktabs}

% ==========================================
% 4. 书籍信息
% ==========================================
\title{\textbf{动机社会：从动机到预测}\\ \large ——社会系统的前向推导与贡献分解（重构版）}
\author{D·declaim}
\date{\today}

% ==========================================
% 5. 正文开始
% ==========================================
\begin{document}

	\maketitle      % 生成封面

	\frontmatter
	\subfile{draft/frontmatter/preface_v2}
	\tableofcontents % 生成目录

	\mainmatter     % 正文页码开始

	% =================================================================
	% 引言：新的目标——预测
	% =================================================================
	\subfile{draft/part1_motivation_goal/introduction_v2}

	% =================================================================
	% PART I: 动机与目标——从行为反推到意向
	% =================================================================
	\part{动机与目标：从行为反推到意向 (Motivation and Goal)}
	\subfile{draft/part1_motivation_goal/chapter1_motivation}
	\subfile{draft/part1_motivation_goal/chapter2_goal}

	% =================================================================
	% PART II: 演化状态空间（P2 → P3 是硬依赖：条件概率必须定义在状态空间上）
	% =================================================================
	\part{演化状态空间：物质、人、相互作用 (The State Space)}
	\subfile{draft/part2_state_space/chapter3_state_triple}
	\subfile{draft/part2_state_space/chapter4_survival}
	\subfile{draft/part2_state_space/chapter5_evolution}

	% =================================================================
	% PART III: 条件概率与贡献分解（原书断裂点的修复处）
	% =================================================================
	\part{条件概率与贡献分解 (Conditional Probability and Contribution)}
	\subfile{draft/part3_conditional_probability/chapter6_prediction_math}
	\subfile{draft/part3_conditional_probability/chapter7_contribution}

	% =================================================================
	% PART IV: 四个通道（嵌套而非并列：每章是上一章「加入一个约束」）
	% =================================================================
	\part{四个通道：模因、经济、政治、法律 (Four Channels)}
	\subfile{draft/part4_channels/chapter8_channel_generator}
	\subfile{draft/part4_channels/chapter9_meme}
	\subfile{draft/part4_channels/chapter10_economy}
	\subfile{draft/part4_channels/chapter11_politics}
	\subfile{draft/part4_channels/chapter12_law}

	% =================================================================
	% PART V: 预测与循环（从贡献表到未来走向）
	% =================================================================
	\part{预测与循环：从贡献表到未来走向 (Forecast and Loop)}
	\subfile{draft/part5_backtest/chapter13_forecast}
	\subfile{draft/part5_backtest/chapter14_closure}

	% =================================================================
	% 数学附录（v2）
	% =================================================================
	\appendix
	\subfile{draft/appendix/appendix_v2_math}

	\backmatter

\end{document}
